\chapter{CRONOGRAMA}

O desenvolvimento deste Trabalho de Conclusão de Curso foi estruturado e executado ao longo do primeiro semestre de 2026. As atividades iniciaram-se de maneira concomitante ao início do semestre letivo e da disciplina de Projeto de Pesquisa Científica e Tecnológica. A Tabela \ref{tab:cronograma} consolida os principais marcos temporais e as fases de evolução do projeto arquitetural do qDATA, desde o delineamento inicial até a sua entrega e finalização no dia atual.

\begin{table}[htb]
\centering
\caption{Cronograma de execução das atividades do TCC}
\label{tab:cronograma}
\begin{tabular}{|p{3cm}|p{12cm}|}
\hline
\textbf{Período} & \textbf{Atividade} \\
\hline
05 de Janeiro & Início do semestre letivo, da disciplina de Projeto de Pesquisa Científica e Tecnológica e definição do escopo macro do LTI FinHub. \\
\hline
Janeiro & Revisão bibliográfica focada em arquitetura de microsserviços, bancos de dados de séries temporais e protocolos de comunicação. \\
\hline
Fevereiro & Preparação da infraestrutura com Docker, conteinerização de dependências. \\
\hline
Março & Desenvolvimento da API assíncrona (REST) utilizando FastAPI e integração com o banco de dados TimescaleDB. \\
\hline
Abril & Implementação da arquitetura gRPC para streaming de dados bidirecional, e configuração de autenticação via JWT/Supabase. \\
\hline
Maio & Testes de carga, validação da idempotência (\textit{Upsert}), e início da escrita dos capítulos textuais (Introdução e Fundamentação). \\
\hline
Junho & Finalização da redação da Metodologia, revisão dos textos produzidos e enquadramento nas normas da ABNT. \\
\hline
22 de Junho & Revisão final, encerramento do desenvolvimento e entrega formal do documento do TCC. \\
\hline
\end{tabular}
\end{table}
